\section{Binding the selected local course}
\label{sec:ex-lattice-selected-course-binding}

The public construction returns a run together with a continuation on its
selected history.  Specializing that continuation to the run's own bolt,
schedule, transfer family and doors yields a binding on the very family used
by the run.  This section proves that same-family assertion and records the
monotonicity needed when comparing paid schedules.  The global schedule and
same-history identities are those of Sections~\ref{sec:common-funded-history}
and~\ref{prop:ex-original-history-attachment}.

Write $R$ for the run selected by the public construction and put
\[
 C_R=\bigl(\operatorname{trail}(T_R),t_R,S_R\bigr),
 \qquad A_R=\mathcal L(C_R),
\]
where $T_R$ is its bolt, $t_R$ its schedule, $S_R$ its transfer family, and
$\mathcal L(C)$ denotes the directed level schedule of a course $C$.  The
family $C_R$ is exactly the family stored in $R$, with its fixed holonomy,
label frames and local sections.

\begin{proposition}[Binding on the selected family]
\label{prop:ex-lattice-selected-course-binding}
The continuation stored in $R$, specialized to its own bolt, schedule,
transfers and doors, supplies a binding $B_R$ for $C_R$.  Its price map is
the selected run's price map.
\end{proposition}

\begin{proof}
The continuation is universally applicable to a bolt, its nested factor,
a schedule, transfer data and doors, subject to the making, climbing and
freight conditions.  The run stores these very objects and receipts: its
tuning supplies nesting, its acceptance data supply the openings, its making
data define the family, and its carrying and freight records supply the
remaining steps.  Specializing all arguments to those stored objects gives
the family $C_R$ by its definition.  The resulting calibrated package
contains a binding $B_R$ on $C_R$, with price map equal to the selected
run's price map.  The binding is constructed from the shape-port, attached
path and shell comparison for that same alloy and bolt, so no separately
selected family is involved.  In assembling the public run, the directed
levels of this family are identified with the levels of the directed alloy;
therefore the schedule used by the binding is $A_R$.
\end{proof}

We record the exact elementary monotonicity for the paid-schedule condition.
For waits $U,V$ over the same bar, write $U\preceq V$ when their requests
obey, pointwise in stage $j$ and scale $s$,
\[
 r_U(j)\le r_V(j),\qquad w_U(j,s)\le w_V(j,s),
 \qquad p_U\le p_V.
\]
Here $r$ is the ready-time request, $w$ the window request, and $p$ the
minimum-period request.  A schedule $A$ has stage map $J_A$, scale values
$s_A(k)$ and period values $P_A(i)$.  Write $\mathsf{Pay}(U,A)$ when the
fixed kernel, bar, stage, initial-stage, length and protection identities
and the entry, scale and volume inequalities hold, together with the
following two conditions:
\begin{enumerate}
\item the separate period condition $p_U\le P_A(i)$ for every $i$;
\item for every $j$ from the required initial stage onward, there is a $k$
      such that
      \begin{equation}\label{eq:ex-lattice-paid-plateau}
      \begin{gathered}
      r_U(j)\le s_A(k),\qquad w_U\bigl(j,s_A(k)\bigr)\le k,\\
      J_A(i)=j\quad
      (k-w_U(j,s_A(k))\le i\le k+w_U(j,s_A(k))).
      \end{gathered}
      \end{equation}
\end{enumerate}
The lower endpoint in this interval uses truncated natural subtraction.
These are the schedule conditions used in
Section~\ref{sec:common-scale-selection}.

\begin{lemma}[Weakening the paid-schedule requests]
\label{lem:ex-lattice-paid-weakening}
If $U\preceq V$ and $\mathsf{Pay}(V,A)$, then $\mathsf{Pay}(U,A)$.
\end{lemma}

\begin{proof}
The fixed identities and entry, scale and volume inequalities do not depend
on the wait.  The separate period condition for $U$ follows for every $i$
from $p_U\le p_V\le P_A(i)$.  Fix a stage $j$ and take the same $k$ supplied
by the plateau condition for $V$.  Its ready-time inequality implies
$r_U(j)\le r_V(j)\le s_A(k)$.  Set
$a=w_U(j,s_A(k))$ and $b=w_V(j,s_A(k))$.  Then $a\le b$, and truncated
subtraction gives
\[
 \max\{0,k-b\}\le \max\{0,k-a\},
 \qquad k+a\le k+b.
\]
Hence the interval of indices on which $J_A(i)=j$ for $V$ contains the
interval required for $U$.  Also $a\le b\le k$, so the window bound for $U$
holds.  Thus the same index $k$ verifies the entire plateau condition for
$U$; this proves $\mathsf{Pay}(U,A)$.
\end{proof}

These results identify the binding family and compare paid-schedule
requirements.  They do not themselves supply the subsequent local-service
estimates.
